\documentclass[10pt,a4paper]{amsart}
\usepackage[margin=19mm]{geometry}
\usepackage{amsmath,amssymb,mathtools}
\usepackage[scr=rsfs,cal=euler]{mathalfa}
\usepackage{xfrac}
\usepackage[unicode,hidelinks,bookmarks=false]{hyperref}
\newcommand{\ie}{i.e.\ }
\newcommand{\cf}{cf.\ }
\newcommand{\resp}{respectively\ }
\newcommand{\Subsection}{\subsection}
\providecommand{\nobreakdash}{\nobreak\hbox{-}}
\setlength{\emergencystretch}{2em}
\multlinegap0pt
\usepackage{amssymb}
\newtheorem{theorem}{Theorem}
\newtheorem{proposition}{Proposition}
\def \T{{\mathbb T}}
\def \Z{{\mathbb Z}}
\title{Localization operators on discrete Orlicz modulation spaces}
\author{Aparajita Dasgupta, Anirudha Poria}
\date{Source: arXiv:2409.05373v2 (CC BY 4.0)}
\begin{document}
\maketitle

\noindent\emph{Source and licence.} Localization operators on discrete Orlicz modulation spaces. Version arXiv:2409.05373v2 (CC BY 4.0), distributed by arXiv under the Creative Commons Attribution 4.0 International licence, http://creativecommons.org/licenses/by/4.0/, as recorded in the arXiv licence metadata for this exact version and shown on the abstract page https://arxiv.org/abs/2409.05373v2. Full licence notice: \texttt{licenses/arxiv-2409.05373-CC-BY-4.0.txt}.\par\smallskip

\noindent\textit{Editorial scope.} Counted unit: the article's proposition pro3 (source0.tex lines 570--573). The article's complete original proof of it is delivered with it (source0.tex lines 574--598, one \texttt{proof} environment of 25 lines). No mathematics is written, repaired or re-derived: the statement and the proof are the article's own lines, in the article's own notation, and no retained line is substituted or re-flowed.  Retained uncounted as the source-local context the proof needs: lines 532--534.  Nothing was omitted from any retained span, so the package carries no omission record.  No retained line contains a citation, so the wrapper carries no bibliography and no bibliography entry.  No retained line contains a figure environment, an image-inclusion command or drawing code, so no image is delivered and the package declares no asset dependency.  Editorial formatting only, in addition: an AMS \texttt{amsart} preamble that restates the article's own declarations and macros used by the retained lines, namely the environments \texttt{theorem}, \texttt{proposition} on separate counters, together with the standard packages those lines need.  The retained environments print in this document as Theorem 1, Proposition 1 in order of appearance, whereas the article prints the same items with the numbering of its own full document.  The complete native source of the article as distributed in the arXiv e-print for this exact version is bundled unchanged as \texttt{sources/164969/source0.tex}.
\begin{theorem}\label{th2}
Let $1 \leq p \leq 2$, $\sigma \in M^p(\Z^n \times \T^n)$, $\Phi$ be a Young function, and $g_1, g_2 \in  M^{\Phi}(\Z^n)$. Then, the operator $\mathfrak{L}^{g_1, g_2}_{\sigma}: \ell^2(\Z^n) \to \ell^2(\Z^n) $ is compact. 
\end{theorem}  

\begin{proposition}\label{pro3}
Let $\sigma \in  M^1(\Z^n \times \T^n)$, $\Phi$ be a Young function, and $g \in  M^{\Phi}(\Z^n)$. Then, the operator $\mathfrak{L}^{g, g}_{\sigma}: \ell^2(\Z^n) \to \ell^2(\Z^n) $ is in $S_1$ and 
\[ \Vert \mathfrak{L}^{g, g}_{\sigma} \Vert_{S_1} \leq 4 \Vert \sigma \Vert_{M^1(\Z^n \times \T^n)} \;  \Vert g \Vert^2_{M^{\Phi}(\Z^n)}. \]
\end{proposition}

\begin{proof}
Let $\sigma \in  M^1(\Z^n \times \T^n)$. Then, from Theorem \ref{th2}, the operator $\mathfrak{L}^{g, g}_{\sigma} \in S_1$. Let $\sigma \in M^1(\Z^n \times \T^n)$ be non-negative real-valued. Then $(({\mathfrak{L}^{g, g}_{\sigma}})^* \mathfrak{L}^{g, g}_{\sigma})^{1/2}=\mathfrak{L}^{g, g}_{\sigma}$. Let $\{v_n\}_n$ be an orthonormal basis for $\ell^2(\Z^n)$ consisting of eigenvalues of the operator $(({\mathfrak{L}^{g, g}_{\sigma}})^* \mathfrak{L}^{g, g}_{\sigma})^{1/2}: \ell^2(\Z^n) \to \ell^2(\Z^n)$. If $\sigma$ is non-negative real-valued function, then by using a similar method as in the proof of Theorem \ref{th2}, we get the following estimate
\begin{eqnarray}\label{eq43}
\Vert \mathfrak{L}^{g, g}_{\sigma} \Vert_{S_1}
& =& \sum_{n=1}^\infty \left\langle (({\mathfrak{L}^{g, g}_{\sigma}})^* \mathfrak{L}^{g, g}_{\sigma})^{1/2}v_n, v_n \right\rangle_{\ell^2(\Z^n)} \nonumber \\
&=& \sum_{n=1}^\infty \left\langle \mathfrak{L}^{g, g}_{\sigma} v_n, v_n \right\rangle_{\ell^2(\Z^n)}  
\leq  \Vert \sigma \Vert_{M^1(\Z^n \times \T^n)} \; \Vert g \Vert^2_{M^{\Phi}(\Z^n)} .
\end{eqnarray}
Next, let $\sigma \in M^1(\Z^n \times \T^n)$ be an arbitrary real-valued function. We write $\sigma = \sigma_+ - \sigma_-$, where $\sigma_+=\max(\sigma, 0)$ and $\sigma_-=-\min(\sigma, 0)$. Then, applying the relation (\ref{eq43}), we get
\begin{eqnarray}\label{eq44}
\Vert \mathfrak{L}^{g, g}_{\sigma} \Vert_{S_1}
& = & \Vert \mathfrak{L}^{g, g}_{\sigma_+} -\mathfrak{L}^{g, g}_{\sigma_-} \Vert_{S_1} \nonumber \\
& \leq & \Vert \mathfrak{L}^{g, g}_{\sigma_+} \Vert_{S_1} +\Vert \mathfrak{L}^{g, g}_{\sigma_-} \Vert_{S_1} \nonumber \\
& \leq &  \Vert g \Vert^2_{M^{\Phi}(\Z^n)} (\Vert \sigma_+ \Vert_{M^1(\Z^n \times \T^n)} + \Vert \sigma_- \Vert_{M^1(\Z^n \times \T^n)}) \nonumber \\
& \leq & 2  \Vert g \Vert^2_{M^{\Phi}(\Z^n)} \Vert \sigma \Vert_{M^1(\Z^n \times \T^n)}.
\end{eqnarray}
Finally, let $\sigma \in  M^1(\Z^n \times \T^n)$ be a complex-valued function. Then, we write $\sigma=\sigma_1+ i \sigma_2$, where $\sigma_1, \sigma_2$ are the real and imaginary parts of $\sigma$ respectively. Applying the relation (\ref{eq44}), we get
\begin{eqnarray*}
\Vert \mathfrak{L}^{g, g}_{\sigma} \Vert_{S_1}
& = & \Vert \mathfrak{L}^{g, g}_{\sigma_1}+i \; \mathfrak{L}^{g, g}_{\sigma_2} \Vert_{S_1}  \\
& \leq & \Vert \mathfrak{L}^{g, g}_{\sigma_1} \Vert_{S_1} + \Vert \mathfrak{L}^{g, g}_{\sigma_2} \Vert_{S_1} \\
&\leq & 2  \Vert g \Vert^2_{M^{\Phi}(\Z^n)} (\Vert \sigma_1 \Vert_{M^1(\Z^n \times \T^n)} + \Vert \sigma_2 \Vert_{M^1(\Z^n \times \T^n)}) \nonumber \\
& \leq & 4 \Vert \sigma \Vert_{M^1(\Z^n \times \T^n)}  \Vert g \Vert^2_{M^{\Phi}(\Z^n)}  .
\end{eqnarray*}
\end{proof}

\end{document}
